\chapter{Security and Reliability}
\label{ch:security}

\section{Security boundary}

A FreeBSD jail is operating-system-level isolation, not hardware virtualization.
It shares the host kernel. Blackship reduces configuration error and limits
resource/network exposure, but cannot strengthen a jail beyond the guarantees
of the running FreeBSD kernel and its policy.

Use a VM boundary for untrusted tenants that require independent kernels or for
workloads whose compromise impact cannot include the host kernel. Keep the host
patched and avoid granting broad jail parameters merely to make an application
work.

\section{Least privilege}

Most inspection and authoring operations are unprivileged. Host mutations need
root, and so do reads that enter a jail or its state: \command{health},
\command{logs}, and \command{gc} (even with \command{--dry-run}). An
unprivileged \command{eva --dry-run} cannot read VNET addresses or kernel
rctl rules and reports them as unverified. The CLI's escalation order is mdo, doas, then sudo, using canonical paths.
The library never escalates.

For interactive administration, define narrowly scoped \command{mac\_do} or
doas policy rather than unrestricted root shells. For a management UI, separate
the unprivileged planning process from the small authorized apply path and pass
\command{--no-escalate} to the former.

\section{Environment and path hygiene}

On mdo re-exec, Blackship resets \command{HOME}, \command{USER}, and
\command{LOGNAME} to root and removes XDG directory overrides. This prevents a
privileged invocation from continuing to use attacker-controlled user paths.

The executable and escalators should be installed in root-owned locations.
Automation should use an absolute Blackship path, a known working directory,
and an explicit manifest. Do not accept an arbitrary manifest path from an
untrusted caller and then authorize its lifecycle operations.

Jailfile copies are bounded to the build context and jail root. Symlink and
parent-directory escapes are rejected. MCP writes use a temporary file in the
target directory followed by rename, preventing an interrupted write from
truncating an existing manifest.

\section{Network policy}

VNET offers a separate network stack; shared networking does not. PF rules
still belong to the host security policy. Audit the base ruleset, NAT, Blackship
anchors, and external interface exposure together.

Port forwarding defaults should bind only where required. A service intended
for localhost should not be published on every host address. DNS and gateway
values in a manifest are configuration, not authentication.

Blackship verifies interface ownership before cleanup and serializes address
leases, reducing races and accidental deletion. Operators should nevertheless
avoid manually renaming or reusing managed interfaces while jails are active.

\section{Storage and secrets}

Jail roots, snapshots, exports, build contexts, staged manifests, and ZFS send
streams may contain secrets. Protect them according to the most sensitive jail
content. A snapshot is not sanitized because the live file was later deleted.

Keep telemetry tokens and local credentials in protected overrides. Review
Jailfiles before building: \command{RUN} is privileged application installation
inside a root filesystem and may fetch or execute external content.

Library-container nullfs mounts deserve two-sided review. Read-write mounts are
\command{nosuid} inside the jail by default and can also be made
\command{noexec}; neither property changes access through the original host
path. Put writable, user-reachable host directories on a dataset with
\command{setuid=off}. Treat \command{allow\_setuid()} and
\command{allow\_chflags()} as narrow exceptions for trusted build workloads,
not general compatibility switches.

\section{Reliability model}

Blackship records ownership before mutation and performs reverse-order rollback.
Cleanup is designed to be idempotent, because a crash can interrupt either the
forward path or rollback itself. Mutable TOML state is written to a unique
same-directory temporary file and atomically renamed over the target. Scope
records are current work-lists; audit records are durable history.

The project treats isolation failures as fatal. It does not continue after a
requested VNET attach or resource policy fails. This fail-closed behavior may
surface host misconfiguration earlier than a hand-written jail script, which is
the intended result.

\section{Safe change procedure}

\begin{enumerate}
  \item identify the exact project, target selector, and resolved full jail names;
  \item validate the manifest and merged overrides;
  \item inspect capabilities and current status;
  \item snapshot or export state that cannot be recreated;
  \item generate a dry-run plan;
  \item apply with non-interactive escalation only when credentials are already
    available;
  \item verify health and network reachability; and
  \item inspect audit and a GC dry run for residual resources.
\end{enumerate}

\section{Backup and recovery}

Test restoration, not only backup creation. A complete recovery plan includes:

\begin{itemize}
  \item the manifest and Jailfiles in version control;
  \item portable exports or replicated ZFS snapshots for mutable jail data;
  \item the pool/dataset naming expected on the destination;
  \item host PF, loader, rc, and escalation policy outside Blackship's data tree;
  \item release availability compatible with the destination kernel; and
  \item a procedure to recreate named networks and forwards with
    \command{blackship setup}.
\end{itemize}

\begin{shipwarning}[Force means proof has been relaxed]
Forced cleanup and garbage collection exist for damaged states where ownership
cannot be established normally. They are not routine shortcuts. Capture status,
scope/audit files, mounts, interfaces, PF anchors, and datasets before forcing;
that evidence is often what explains the incident.
\end{shipwarning}

\section{Auditability}

Use stable JSON error codes and audit events for automated decisions. Retain
the package version and capability report alongside incident logs. Human CLI
text is for diagnosis and may evolve; it is not a durable event schema.
